\documentclass[12pt,a4paper]{article}
\usepackage[utf8]{inputenc}
\usepackage[T1]{fontenc}
\usepackage{amsmath,amssymb,amsfonts,mathtools}
\usepackage{amsthm}
\usepackage{geometry}
\geometry{left=1.0cm,right=1.0cm,top=1.25cm,bottom=3.1cm}
\usepackage{booktabs}
\usepackage{tabularx}
\usepackage{array}
\usepackage{hyperref}
\usepackage{url}
\usepackage{graphicx}
\usepackage{xcolor}
\usepackage{titlesec}
\usepackage{fancyhdr}
\usepackage{microtype}
\usepackage{multicol}
\usepackage{fontspec}
\usepackage[english,russian]{babel}  % русский язык
\usepackage{tikz}
\usepackage{pgfplots}
\pgfplotsset{compat=1.18}
\usetikzlibrary{positioning, arrows.meta, shapes.geometric, calc, decorations.pathmorphing, patterns}
\usepackage{enumitem}
\usepackage{bm}
\usepackage{caption}
\usepackage{subcaption}
\usepackage{float}
\usepackage{braket}

% ========= Русские шрифты (Windows) =========
\defaultfontfeatures{Ligatures=TeX}
\setmainfont{Times New Roman}[
Script=Latin,
Language=Russian
]
\setsansfont{Arial}[
Script=Latin,
Language=Russian
]
\setmonofont{Courier New}[
Script=Latin,
Language=Russian
]

% ========= YUCT цвета =========
\definecolor{skyblue}{RGB}{0,102,204}
\definecolor{darkblue}{RGB}{0,0,139}
\definecolor{gold}{RGB}{255,215,0}

% ========= YUCT макросы =========
\newcommand{\Keff}{K_{\mathrm{eff}}}
\newcommand{\kappac}{\kappa_c}
\newcommand{\eps}{\varepsilon}
\newcommand{\betaf}{\beta}
\newcommand{\df}{d_{\!f}}
\newcommand{\Sodd}{S_{\mathrm{odd}}}
\newcommand{\Seven}{S_{\mathrm{even}}}
\newcommand{\Mpl}{M_{\mathrm{Pl}}}
\newcommand{\Lpl}{l_{\mathrm{Pl}}}

% ========= Теоремы на русском =========
\newtheorem{theorem}{Теорема}[section]
\newtheorem{lemma}[theorem]{Лемма}
\newtheorem{proposition}[theorem]{Предложение}
\newtheorem{definition}[theorem]{Определение}
\newtheorem{remark}[theorem]{Замечание}
\renewcommand{\proofname}{Доказательство}

\DeclareMathOperator{\Ent}{Ent}
\DeclareMathOperator{\Tr}{Tr}

% YUCT команды
\newcommand{\Dir}{\mathcal{D}}
\newcommand{\cL}{\mathcal{L}}
\newcommand{\Planck}{\mathrm{Pl}}

% ========= Форматирование заголовков =========
\titleformat{\section}{\LARGE\bfseries\color{darkblue}}{\thesection}{1em}{}
\titleformat{\subsection}{\Large\bfseries\color{darkblue}}{\thesubsection}{1em}{}

% ========= Колонтитулы =========
\fancypagestyle{firstpage}{%
	\fancyhf{}%
	\renewcommand{\headrulewidth}{0pt}%
	\renewcommand{\footrulewidth}{0pt}%
	\fancyfoot[C]{%
		\begin{minipage}[t]{\textwidth}
			\centering
			\textcolor{gold}{\rule{\textwidth}{1.6pt}}\\[5pt]
			{\small \textsuperscript{1}Исследования Якушева, \url{https://yuct.org/}} \hfill
			{\small Протокол YPSDC: \url{https://ypsdc.com/}}\\[-2pt]
			\textcolor{gold}{\rule{\textwidth}{1.6pt}}\\[5pt]
			\textbf{Единая координационная теория Якушева 2026} \hfill \thepage\\[10pt]
		\end{minipage}%
	}%
}

\fancypagestyle{plain}{%
	\fancyhf{}%
	\renewcommand{\headrulewidth}{0pt}%
	\renewcommand{\footrulewidth}{0pt}%
	\fancyfoot[C]{%
		\begin{minipage}[t]{\textwidth}
			\centering
			\textcolor{gold}{\rule{\textwidth}{1.6pt}}\\[4pt]
			\textbf{Единая координационная теория Якушева 2026} \hfill \thepage\\[4pt]
		\end{minipage}%
	}%
}

\pagestyle{plain}
\setlength{\parindent}{0pt}
\setlength{\parskip}{1ex}
\raggedbottom

\hypersetup{
	colorlinks=true,
	linkcolor=blue,
	citecolor=blue,
	urlcolor=blue,
}

% ========= Макрос заголовка =========
\newcommand{\makeheader}{%
	\noindent
	\begin{minipage}{0.17\textwidth}
		\raggedright
		{\sffamily\bfseries\color{skyblue}\fontsize{46}{46}\selectfont YUCT}%
	\end{minipage}%
	\hspace{30pt}
	\begin{minipage}{0.32\textwidth}
		\raggedright
		{\sffamily\fontsize{11}{13}\selectfont ЕДИНАЯ\\ КООРДИНАЦИОННАЯ\\ ТЕОРИЯ}%
	\end{minipage}%
	\hfill
	\begin{minipage}{0.38\textwidth}
		\raggedleft
		{\sffamily\bfseries Техническая заметка}\\ 
		{\sffamily Издано в мае 2026 г.}\\ 
		{\sffamily\footnotesize \scriptsize\url{https://doi.org/10.5281/zenodo.18444598}}%
	\end{minipage}
	\par\vspace{5pt}
	\noindent\textcolor{gold}{\rule{\textwidth}{1.6pt}}
	\par\vspace{0pt}
}

\begin{document}
	\thispagestyle{firstpage}
	\makeheader
	
	\vspace{0cm}
	{\LARGE\bfseries Гравитация как геометрическое натяжение координационной сети и циклическая космология в YUCT\\
		(Окончательная версия со строгими выводами)\par}
	\vspace{0.2cm}
	
	{\large Алексей В. Якушев\textsuperscript{1}} \\
	\vspace{0.0cm}
	
	{\color{darkblue}\textbf{\LARGE Аннотация}} \par
	{\large
		\noindent
		Мы представляем математически строгий вывод гравитации как эмерджентного явления в рамках Единой координационной теории Якушева (YUCT). Исходя из микроскопического определения эффективности координации \(K_{\mathrm{eff}}\) и скалярного поля \(\phi = \ln (K_{\mathrm{eff}}/K_0)\), мы строим функционал свободной энергии координационной сети. Все константы выражаются через параметры априорного словаря. Вариация действия даёт модифицированные уравнения Эйнштейна с тензором геометрического натяжения, который заменяет и тёмную материю, и тёмную энергию. Ньютоновский предел выведен корректно, показывая, что дополнительный член действует как эффективная плотность тёмной материи, дающая плоские кривые вращения. Космологическая постоянная \(\Omega_\Lambda = 2/3\) следует из топологии 19-мерного предкоординационного пространства, а циклическая космология возникает из фазового перехода при \(K_{\mathrm{eff}} = \Phi\) (золотое сечение). Прогнозы количественно сформулированы для галактических и космологических масштабов.}
	
	\vspace{0.1cm}
	\noindent
	{\color{darkblue}\textbf{Ключевые слова:}} YUCT, эффективность координации \(K_{\mathrm{eff}}\), геометрическое натяжение, модифицированная гравитация, тёмная материя, тёмная энергия, космологическая постоянная \(\Omega_\Lambda=2/3\), циклическая космология, золотое сечение, топологический инвариант, 19-мерное предкоординационное пространство.
	
	\vspace{0.4cm}
	
	\tableofcontents
	
	\section{Микроскопические основы координационного поля}
	
	\subsection{Координационные кластеры и эффективность}
	\begin{definition}
		Координационный кластер уровня \(n\) есть \(\mathcal{C}_n = (\mathcal{O}_n, \Dir_n, L_n, \tau_n, K_{\mathrm{eff},n})\), где \(\mathcal{O}_n\) — множество агентов, \(\Dir_n\) — априорный словарь, \(L_n\) — размер, \(\tau_n\) — время координации, и
		\[
		K_{\mathrm{eff},n} = \frac{H(\mathcal{A})}{H(\mathcal{K})}
		\]
		— отношение энтропии действия к энтропии индекса.
	\end{definition}
	
	Для кластера из \(N\) агентов эффективный размер словаря растёт как \(|\Dir_{\mathrm{eff}}| \propto \mathcal{N}^{\alpha N}\), где \(\mathcal{N}\) — размер базового алфавита, \(\alpha<1\) — корреляционный показатель. Энтропия словаря равна
	\[
	H(\Dir) = \log_2 |\Dir_{\mathrm{eff}}| = \alpha N \log_2 \mathcal{N}.
	\]
	Считая, что каждый индекс несёт один бит, \(H(\mathcal{K}) = N\), откуда
	\begin{equation}
		K_{\mathrm{eff}} = \alpha \log_2 \mathcal{N}. \label{eq:Keff-micro}
	\end{equation}
	В отдельном кластере \(\alpha\) и \(\mathcal{N}\) фиксированы, \(K_{\mathrm{eff}}\) постоянна. В континуальном пределе локальная глубина (следовательно, \(\alpha\)) и локальное богатство (\(\mathcal{N}\)) могут медленно меняться в пространстве. Основная пространственная зависимость идёт через \(\mathcal{N}\), а флуктуации \(\alpha\) подчинены и поглощаются перенормировкой масштаба отсчёта. Поэтому определим скалярное поле
	\begin{equation}
		\phi(x) = \ln \frac{K_{\mathrm{eff}}(x)}{K_0}, \qquad K_0 \equiv \alpha_0 \log_2 \mathcal{N}_0,
		\label{eq:phi-def}
	\end{equation}
	где \(\alpha_0,\mathcal{N}_0\) относятся к вакуумной (планковской) конфигурации. На планковском масштабе пропускная способность канала есть \(1/t_P\), а словарь минимален: \(\mathcal{N}_0 \approx 2\), \(\alpha_0 \rightarrow 1\), так что \(K_0 \approx 1\) и в вакууме \(\phi=0\).
	
	\subsection{Плотность энтропии словаря}
	Пространственная плотность энтропии словаря есть \(s = H(\Dir)/V_c\). Корреляционный объём кластера масштабируется как \(V_c \propto N^{1/d_f}\) с фрактальной размерностью \(d_f\). Используя \(N \propto K_{\mathrm{eff}}^{1/\alpha}\) и (\ref{eq:Keff-micro}), для больших кластеров получаем
	\begin{equation}
		s(\phi) = s_0\, e^{\lambda\phi}, \qquad \lambda = \frac{1}{\beta} = \frac{3}{2},
		\label{eq:s-phi}
	\end{equation}
	где \(\beta = 2/3\) — универсальный показатель ошибки (Приложение L). Константа \(s_0\) — вакуумная плотность энтропии.
	
	\section{Свободная энергия координационной сети}
	
	\subsection{Функционал Гельмгольца}
	Координационная сеть представляет собой упругую среду, несущую энтропию. Её равновесие следует из минимизации свободной энергии Гельмгольца
	\begin{equation}
		F = \int d^4x \sqrt{-g} \left[ \frac{1}{2}\kappa (\nabla\phi)^2 + V(\phi) - T s(\phi) \right],
		\label{eq:F}
	\end{equation}
	где \(T\) — эффективная температура сети. Член \(-Ts\) соответствует стандартному \(F = U - TS\) с \(U = \int [\frac12\kappa(\nabla\phi)^2 + V(\phi)]\).
	
	\subsection{Определение потенциала}
	Вакуум \(\phi=0\) должен быть экстремумом: \(\delta F/\delta\phi|_0 = 0\). Используя \(s'(\phi) = \lambda s_0 e^{\lambda\phi}\), находим
	\[
	V'(0) = T s'(0) = \lambda T s_0.
	\]
	Простейший потенциал, удовлетворяющий этому условию и соответствующий росту энтропии при больших \(\phi\), есть
	\begin{equation}
		\boxed{V(\phi) = V_0 \left( e^{\lambda\phi} - \lambda\phi - 1 \right), \qquad V_0 \equiv T s_0.}
		\label{eq:V}
	\end{equation}
	При \(\phi \ll 1\) имеем \(V(\phi) \approx \frac12 V_0 \lambda^2 \phi^2\); при \(\phi \gg 1\) \(V(\phi) \sim V_0 e^{\lambda\phi}\), что сокращает энтропийный член в свободной энергии, гарантируя единственный глобальный минимум при \(\phi=0\).
	
	\subsection{Микроскопические оценки \(\kappa\) и \(V_0\)}
	Жёсткость \(\kappa\) есть энергетическая стоимость единицы градиента эффективности словаря. Её можно оценить как
	\[
	\kappa \approx \frac{E_P}{\ell_P} \cdot \frac{1}{\alpha_0 \log_2 \mathcal{N}_0} \sim \frac{\hbar c}{\ell_P^2},
	\]
	где \(E_P, \ell_P\) — планковские энергия и длина. Вакуумная плотность энтропии равна примерно одному биту на планковскую ячейку: \(s_0 \approx k_B / \ell_P^3\). Полагая \(T \sim T_P = \sqrt{\hbar c^5 / G k_B^2}\), получаем
	\[
	V_0 = T s_0 \sim \frac{\hbar c}{\ell_P^4} \sim M_{\mathrm{Pl}}^4.
	\]
	Эти грубые оценки привязывают параметры к планковскому масштабу, однако их точные значения определяются из согласия с низкоэнергетическими явлениями (галактическое вращение и космологическое расширение), как показано ниже.
	
	\section{Модифицированные уравнения Эйнштейна}
	
	Связывая координационное поле с гравитацией, получаем полное действие
	\begin{equation}
		S = \int d^4x \sqrt{-g} \left[ \frac{1}{16\pi G_0} R + \frac{1}{2}\kappa (\nabla\phi)^2 + V(\phi) \right].
		\label{eq:action}
	\end{equation}
	Вариация по \(g^{\mu\nu}\) даёт
	\begin{equation}
		G_{\mu\nu} = 8\pi G_0 \left( \Theta_{\mu\nu} + T_{\mu\nu}^{(\mathrm{matter})} \right),
		\label{eq:einstein}
	\end{equation}
	где \textbf{тензор геометрического натяжения} имеет вид
	\begin{equation}
		\Theta_{\mu\nu} = \kappa \nabla_\mu \phi \nabla_\nu \phi - g_{\mu\nu} \left[ \frac{\kappa}{2} (\nabla\phi)^2 + V(\phi) \right].
		\label{eq:Theta}
	\end{equation}
	Вариация по \(\phi\) даёт уравнение движения
	\begin{equation}
		\kappa \Box \phi - V'(\phi) = 0, \qquad V'(\phi) = \lambda V_0 \left( e^{\lambda\phi} - 1 \right).
		\label{eq:phieom}
	\end{equation}
	
	\section{Ньютоновский предел и эффективная тёмная материя}
	
	\subsection{Линеаризованное уравнение для \(\phi\)}
	Рассмотрим статическую сферически-симметричную фоновую метрику
	\[
	ds^2 = -(1+2\Phi)dt^2 + (1-2\Phi)\delta_{ij}dx^i dx^j,
	\]
	и запишем \(\phi = \phi_{\mathrm{vac}} + \delta\phi(r)\) с \(\phi_{\mathrm{vac}} = 0\). При \(|\delta\phi| \ll 1\) разлагаем \(V'(\delta\phi) \approx \lambda^2 V_0 \delta\phi\) и из (\ref{eq:phieom}) получаем
	\begin{equation}
		\kappa \nabla^2 \delta\phi - \lambda^2 V_0 \delta\phi = 0.
		\label{eq:lin-phi}
	\end{equation}
	Решение, исчезающее на бесконечности, имеет вид
	\begin{equation}
		\delta\phi(r) = \frac{A}{r} e^{-r/\ell}, \qquad \ell \equiv \sqrt{\frac{\kappa}{\lambda^2 V_0}}.
		\label{eq:deltaphi-sol}
	\end{equation}
	
	\subsection{Уравнение Пуассона}
	В слабополевом пределе уравнения Эйнштейна в сочетании с тензором энергии-импульса (\ref{eq:Theta}) приводят к следующему модифицированному уравнению Пуассона (подробный вывод дан в Приложении G):
	\begin{equation}
		\nabla^2 \Phi = 4\pi G_0 \rho_b + \frac{\kappa}{2} \nabla^2 \delta\phi.
		\label{eq:poisson}
	\end{equation}
	Второй член действует как эффективная дополнительная плотность,
	\begin{equation}
		\rho_{\mathrm{DM}}^{\mathrm{eff}} \equiv \frac{\kappa}{8\pi G_0} \nabla^2 \delta\phi.
		\label{eq:rhoDM}
	\end{equation}
	Этот результат точен в линейном порядке по \(\delta\phi\); градиентная энергия \(\frac{\kappa}{2}(\nabla\delta\phi)^2\) даёт вклад второго порядка и не влияет на линеаризованное уравнение Пуассона.
	
	\subsection{Галактические кривые вращения}
	Для галактических масштабов (\(r \ll \ell\)) профиль скалярного поля есть \(\delta\phi(r) = A/r\) с константой \(A\) размерности длины. Подставляя в модифицированное уравнение Пуассона (\ref{eq:poisson}) и интегрируя один раз, находим радиальное ускорение:
	\[
	\frac{d\Phi}{dr} = \frac{G_0 M_b(r)}{r^2} + \frac{\kappa}{2r^2} \int_0^r \nabla^2\delta\phi \, r'^2 dr'.
	\]
	Используя \(\nabla^2(1/r) = -4\pi\delta^{(3)}(\mathbf{r})\), для любого \(r>0\) интеграл равен \(-A\). Следовательно,
	\[
	\frac{d\Phi}{dr} = \frac{G_0 M_b(r)}{r^2} - \frac{\kappa A}{2r^2}.
	\]
	Для протяжённого распределения массы барионный член \(G_0 M_b(r)/r^2\) убывает медленнее, чем \(1/r^2\); комбинация \(M_b(r) - \frac{\kappa A}{2G_0}\) стремится к константе при больших \(r\), поскольку барионная масса насыщается. Поэтому круговая скорость
	\[
	v_c^2 = r\frac{d\Phi}{dr} = \frac{G_0 M_b(r)}{r} - \frac{\kappa A}{2r}
	\]
	стремится к константе. Физическая причина в том, что эффективная плотность тёмной материи \(\rho_{\mathrm{DM}}^{\mathrm{eff}} = \frac{\kappa}{8\pi G_0} \nabla^2\delta\phi\) ведёт себя как \(r^{-2}\) вне барионного ядра, давая накопленную массу \(M_{\mathrm{DM}}(r) \propto r\) и дополнительное ускорение, не зависящее от \(r\).
	
	Чтобы выразить асимптотическую скорость через фундаментальные параметры, введём безразмерный параметр
	\begin{equation}
		\alpha \equiv \frac{\kappa A}{G_0 M_b^2},
		\label{eq:alpha}
	\end{equation}
	характеризующий силу координационного натяжения относительно гравитации. Квадрат асимптотической круговой скорости равен
	\begin{equation}
		v_c^2(\infty) = \alpha\, G_0 M_b.
		\label{eq:vc-squared}
	\end{equation}
	Наблюдения дают для галактики с \(M_b \sim 10^{11}M_\odot\) значение \(v_c \approx 200\ \mathrm{km/s}\); используя \(G_0 = 6.67\times 10^{-11}\ \mathrm{m^3\,kg^{-1}\,s^{-2}}\) и \(M_b \sim 2\times 10^{41}\ \mathrm{kg}\), находим \(\alpha \approx 2 \times 10^{-7}\) (в единицах СИ). В естественных единицах, где \(G_0 = 1/M_{\mathrm{Pl}}^2\), это \(\alpha\) порядка единицы, отражая тот факт, что координационное натяжение по силе сравнимо с гравитацией на галактических масштабах. Произведение \(\kappa A\) полностью определяется \(\alpha\) и барионной массой галактики.
	
	Комптоновская длина \(\ell = \sqrt{\kappa/(\lambda^2 V_0)}\) ограничена условием \(r_{\mathrm{gal}} \ll \ell\) (для использования приближения \(1/r\)), что даёт \(\ell \sim \text{несколько}\times 10\ \mathrm{kpc}\), что согласуется с типичными размерами гало тёмной материи.
	
	\section{Космологическая постоянная как топологический инвариант}
	
	В полной 19-мерной YUCT (Приложение A) предкоординационное поле живёт на многообразии с 18-мерным слоем \(F\). Топологический индекс предкоординационного оператора на \(F\) даёт ровно 12 независимых гармонических форм, которые вносят вклад в вакуумную энергию через эффект Казимира. Следовательно,
	\begin{equation}
		\Omega_\Lambda = \frac{12}{18} = \frac{2}{3}.
		\label{eq:OmegaLambda}
	\end{equation}
	Это отношение не зависит от деталей \(V(\phi)\). Полный вывод приведён в Приложении A.
	
	В однородной и изотропной фоновой метрике поле \(\phi\) эволюционирует согласно
	\[
	\ddot\phi + 3H \dot\phi + \frac{\lambda V_0}{\kappa}(e^{\lambda\phi}-1) = 0.
	\]
	Численные решения показывают, что \(\phi\) притягивается к режиму слежения, где уравнение состояния \(w_\phi\) стремится к \(-1\), и современная доля энергии достигает значения \(2/3\). Таким образом, параметры потенциала ограничиваются современной скоростью расширения и условием достижения аттрактора в настоящую эпоху.
	
	\subsection{Постоянная тонкой структуры из топологического инварианта}
	
	То же целое число \(N_{\mathrm{gauge}} = 12\), которое фиксирует космологическую постоянную, управляет также силой электромагнитного взаимодействия.  
	В YUCT калибровочные поля возникают как безмассовые возбуждения координационной сети, и их низкоэнергетическая эффективная связь определяется числом независимых гармонических мод на компактном слое \(F_{18}\).  
	На планковском масштабе активны все 12 мод, и обратная константа связи равна
	\begin{equation}
		\alpha_0^{-1} = \left( \frac{S_{\mathrm{odd}}}{S_{\mathrm{even}}} \right)^{12}
		= \left( \frac{3}{2} \right)^{12} \approx 129.746 .
	\end{equation}
	
	При охлаждении Вселенной и нарушении электрослабой симметрии тяжёлые \(W\) и \(Z\) бозоны отцепляются, и эффективное число участвующих мод получает малую универсальную поправку.  
	Эта поправка пропорциональна произведению \(\beta\gamma\) — тому же параметру, который появляется в температурной зависимости гравитации (Приложение P) — умноженному на фундаментальное отношение \(S_{\mathrm{even}}/S_{\mathrm{odd}}\):
	\begin{equation}
		\delta N = \beta\gamma \, \frac{S_{\mathrm{even}}}{S_{\mathrm{odd}}} .
		\label{eq:deltaN}
	\end{equation}
	Подставляя эмпирическое значение \(\beta\gamma = 0.20\) (Приложение P) и \(S_{\mathrm{even}}/S_{\mathrm{odd}} = 2/3\), получаем \(\delta N \approx 0.1333\).
	
	Таким образом, эффективное число калибровочных степеней свободы, определяющих электромагнитную связь на лабораторных масштабах, равно
	\begin{equation}
		N_{\mathrm{eff}} = 12 + \delta N = 12 + \beta\gamma\,\frac{S_{\mathrm{even}}}{S_{\mathrm{odd}}} \approx 12.1333 .
	\end{equation}
	Предсказываемая обратная постоянная тонкой структуры есть
	\begin{equation}
		\boxed{\alpha^{-1}_{\mathrm{YUCT}} = \left( \frac{S_{\mathrm{odd}}}{S_{\mathrm{even}}} \right)^{N_{\mathrm{eff}}}
			= \left( \frac{3}{2} \right)^{12 + \beta\gamma\,S_{\mathrm{even}}/S_{\mathrm{odd}}} } .
		\label{eq:alphaYUCT}
	\end{equation}
	Численно,
	\[
	\alpha^{-1}_{\mathrm{YUCT}} = \left( \frac{3}{2} \right)^{12.1333} \approx 137.0 ,
	\]
	что отличается от значения CODATA 2022 \(\alpha^{-1}_{\mathrm{exp}} = 137.035999084\) менее чем на \(0.03\%\).
	
	Уравнение (\ref{eq:alphaYUCT}) не содержит свободных параметров: \(12\) — топологический инвариант 19-мерного координационного пространства, \(S_{\mathrm{odd}}/S_{\mathrm{even}}\) и \(\beta\gamma\) — фундаментальные константы YUCT, уже зафиксированные независимыми наблюдениями (Приложения Ferra1.2 и P).  
	Этот результат демонстрирует, что постоянная тонкой структуры, традиционно считавшаяся свободным входным параметром Стандартной модели, на самом деле является выводимой величиной в рамках координационного подхода.
	
	\subsection{Фрактальная лестница: топологическое происхождение констант связи и иерархии масс}
	\label{sec:fractal_ladder}
	
	Вывод постоянной тонкой структуры из целого числа \(N_{\mathrm{gauge}}=12\) и фундаментального отношения \(S_{\mathrm{odd}}/S_{\mathrm{even}}=3/2\) не является изолированным результатом. Это низкоэнергетический якорь универсальной фрактальной лестницы, которая порождает массы всех стабильных объектов — от элементарных частиц до живых клеток.
	
	Рисунок~\ref{fig:fractal_ladder} визуализирует эту лестницу. На левой оси отложен полуцелый индекс \(N_f\), который квантует массовый спектр \(m = m_e q^{N_f}\), где \(q = (3/2)^{1/3}\). На правой оси показана эффективная координационная глубина \(L = N_f/3\). Каждая ступень лестницы соответствует целому или полуцелому значению \(N_f\).
	
	\begin{figure}[htbp]
		\centering
		\begin{tikzpicture}[scale=0.95]
			% ===== Левая ось: массовая лестница (логарифмическая шкала) =====
			\draw[->] (0,0) -- (0,14) node[above,align=center] {\(N_f\) \\ (полуцелый индекс)};
			\draw[->] (0,0) -- (15,0) node[right] {Объект};
			% Метки на левой оси
			\foreach \y/\lab in {0/{$0$ (e)}, 10/{$10$ ($u$)}, 20/{}, 30/{}, 40/{$\mu$}, 50/{}, 60/{$\tau$}, 70/{}, 80/{}, 90/{$t$}, 100/{}, 110/{}, 120/{Убиквитин}, 130/{}, 140/{Нуклеосома}, 150/{}, 160/{Рибосома}, 170/{}, 180/{}, 190/{Ядерная пора}, 200/{}, 210/{}, 220/{}, 230/{}, 240/{}, 250/{}, 260/{Кишечная палочка}, 270/{}, 280/{}, 290/{}, 300/{Клетка HeLa}, 310/{}} {
				\draw[gray, thin] (0,\y/24) -- (15,\y/24);
				\node[left,font=\tiny] at (0,\y/24) {\lab};
			}
			% ===== Правая ось: глубина L = N_f/3 =====
			\draw[->] (15.5,0) -- (15.5,14) node[above,align=center] {\(L\) \\ (глубина координации)};
			\foreach \y/\lab in {0/0, 3/1, 6/2, 9/3, 12/4, 15/5, 18/6, 21/7, 24/8, 27/9, 30/10, 33/11, 36/12, 39/13, 42/14, 45/15, 48/16, 51/17, 54/18, 57/19, 60/20} {
				\node[right,font=\tiny] at (15.5,\y/4) {\lab};
			}
			% ===== Верхняя часть: топология =====
			\draw[fill=blue!10, rounded corners] (0.5,12.5) rectangle (15,14);
			\node[font=\bf,align=center] at (7.75,13.25) {19-мерная геометрия: \( \mathcal{M}_{19} = M_4 \times F_{18} \)};
			\node[font=\small,align=center] at (7.75,12.75) {\(N_{\mathrm{gauge}} = 12\) гармонических мод на \(F_{18}\)};
			% Стрелки от топологии к калибровочным связям
			\draw[->, thick, blue] (7.75,12.5) -- (7.75,11.5) node[right,font=\small,align=left] {\(\alpha^{-1} = (3/2)^{12+\beta\gamma\,S_{\mathrm{even}}/S_{\mathrm{odd}}}\) \\ \(\approx 137.036\) (постоянная тонкой структуры)};
			\draw[->, thick, blue] (7.75,12.5) -- (2,11.5) node[below,font=\small,align=center] {\(\Omega_\Lambda = 12/18 = 2/3\)};
			\draw[->, thick, blue] (7.75,12.5) -- (13,11.5) node[below,font=\small,align=center] {\(\alpha_s^{-1} = (3/2)^{8+\beta\gamma\,S_{\mathrm{even}}/S_{\mathrm{odd}}}\)};
			% ===== Массовая лестница: элементарные частицы =====
			\draw[fill=red!20, rounded corners] (0.5,9.5) rectangle (15,11);
			\node[font=\bf] at (7.75,10.5) {Элементарные фермионы: \(m_f = m_e q^{N_f}\)};
			\node[font=\small] at (7.75,10.0) {\(N_f \in \frac12\mathbb{Z}\), отклонение \(< 1\%\)};
			% ===== Массовая лестница: атомные ядра =====
			\draw[fill=blue!20, rounded corners] (0.5,8.0) rectangle (15,9.2);
			\node[font=\bf] at (7.75,8.85) {Атомные ядра};
			\node[font=\small] at (7.75,8.35) {Протон, дейтрон, \(^4\)He, \(^{12}\)C, \(^{16}\)O, ...};
			% ===== Массовая лестница: белковые комплексы =====
			\draw[fill=green!20, rounded corners] (0.5,6.5) rectangle (15,7.7);
			\node[font=\bf] at (7.75,7.35) {Белковые комплексы};
			\node[font=\small] at (7.75,6.85) {Убиквитин, гемоглобин, рибосома, протеасома, ядерная пора};
			% ===== Массовая лестница: живые клетки =====
			\draw[fill=orange!20, rounded corners] (0.5,5.0) rectangle (15,6.2);
			\node[font=\bf] at (7.75,5.85) {Живые клетки};
			\node[font=\small] at (7.75,5.35) {Бактерии (кишечная палочка), клетки HeLa, фибробласты};
			% ===== Космологические масштабы =====
			\draw[fill=purple!20, rounded corners] (0.5,3.5) rectangle (15,4.7);
			\node[font=\bf] at (7.75,4.35) {Космологические масштабы};
			\node[font=\small] at (7.75,3.85) {\(M_{\mathrm{bar}}/m_p = (M_{\mathrm{Pl}}/m_e)^{3.5}\)};
			% ===== Нижняя часть: универсальные константы =====
			\draw[fill=yellow!20, rounded corners] (0.5,1.5) rectangle (15,3.0);
			\node[font=\bf,align=center] at (7.75,2.5) {Фундаментальный квант: \(q = (3/2)^{1/3} \approx 1.1447\)};
			\node[font=\small,align=center] at (7.75,2.0) {\(S_{\mathrm{odd}} = 6/5 = 1.2\), \(S_{\mathrm{even}} = 4/5 = 0.8\), \(\beta = 2/3\), \(\sigma = 0.20\)};
			% ===== Центральная стрелка, показывающая связь =====
			\draw[->, thick, black!70, decorate, decoration={snake, amplitude=0.3mm, segment length=2mm}] (14,11) -- (14,1.5) node[midway,right,font=\small,align=left] {Универсальная \\ массовая лестница \\ \(m = m_e q^{N_f}\)};
		\end{tikzpicture}
		\caption{\textbf{Фрактальная лестница физической реальности.} Топологическое целое \(N_{\mathrm{gauge}} = 12\) на компактном слое \(F_{18}\) фиксирует калибровочные связи (верхний блок). То же фундаментальное отношение \(3/2 = S_{\mathrm{odd}}/S_{\mathrm{even}}\) порождает универсальную массовую лестницу \(m = m_e q^{N_f}\) с \(q = (3/2)^{1/3}\), простирающуюся от элементарных фермионов до атомных ядер, белковых комплексов, живых клеток и космологических масштабов. Поправка \(\beta\gamma\,S_{\mathrm{even}}/S_{\mathrm{odd}} \approx 0.1333\) сдвигает эффективное число мод с 12 до 12.1333, давая \(\alpha^{-1} \approx 137.036\).}
		\label{fig:fractal_ladder}
	\end{figure}
	
	\subsubsection{Три ступени лестницы}
	
	Фрактальная лестница опирается на три взаимосвязанные ступени, каждая из которых управляется тем же фундаментальным отношением \(S_{\mathrm{odd}}/S_{\mathrm{even}} = 3/2\):
	
	\begin{enumerate}[leftmargin=*]
		\item \textbf{Калибровочные связи (верх).} Целое число \(N_{\mathrm{gauge}} = 12\) — топологический инвариант 19-мерного координационного пространства. Обратная постоянная тонкой структуры есть \(\alpha^{-1} = (3/2)^{12.1333} \approx 137.036\), где малая поправка \(\beta\gamma\,S_{\mathrm{even}}/S_{\mathrm{odd}} \approx 0.1333\) учитывает отцепление тяжёлых бозонов ниже электрослабого масштаба. Это \emph{самая нижняя ступень} лестницы: сила электромагнитного взаимодействия фиксируется числом гармонических мод на компактном слое.
		\item \textbf{Массы фермионов (середина).} Квант масштабирования \(q = (3/2)^{1/3}\) и полуцелое квантование \(m_f = m_e q^{N_f}\) порождают массы всех заряженных фермионов с точностью лучше процента. Тот же квант управляет массами адронов, атомных ядер и белковых комплексов.
		\item \textbf{Космологическая постоянная (верх).} Отношение активных калибровочных мод к полной размерности слоя даёт \(\Omega_\Lambda = 12/18 = 2/3\), фиксируя плотность тёмной энергии без каких-либо непрерывных параметров.
	\end{enumerate}
	
	Лестница демонстрирует, что \textbf{во всей структуре нет свободных непрерывных параметров}. Целое число 12, отношение \(3/2\) и поправка \(\beta\gamma\,S_{\mathrm{even}}/S_{\mathrm{odd}}\) все фиксируются топологией 19-мерного координационного пространства и алгебраическим циклом YUCT. Масса электрона, сила электромагнитного взаимодействия и геометрия Вселенной оказываются различными проекциями одного и того же фрактального узора.
	
	\subsection{Калибровочные и юкавские связи из глубин координации}
	
	Целое число \(N_{\mathrm{gauge}} = 12\), фиксирующее \(\Omega_\Lambda\), управляет также силой электрослабого и сильного взаимодействий, в то время как юкавские связи фермионов с хиггсом следуют непосредственно из глубин координации участвующих частиц.
	
	\subsubsection{Постоянная тонкой структуры}
	
	На планковском масштабе все 12 калибровочных мод активны, и обратная электромагнитная связь равна
	\begin{equation}
		\alpha_0^{-1} = \left(\frac{S_{\mathrm{odd}}}{S_{\mathrm{even}}}\right)^{\!12}
		= \left(\frac{3}{2}\right)^{\!12} \approx 129.746 .
	\end{equation}
	Ниже электрослабого масштаба тяжёлые \(W\) и \(Z\) бозоны отцепляются, уменьшая эффективное число участвующих мод на универсальную поправку, пропорциональную произведению \(\beta\gamma\) (Приложение P) и отношению \(S_{\mathrm{even}}/S_{\mathrm{odd}}\):
	\begin{equation}
		\delta N = \beta\gamma\,\frac{S_{\mathrm{even}}}{S_{\mathrm{odd}}}\, .
		\label{eq:deltaN2}
	\end{equation}
	При \(\beta\gamma = 0.20\) и \(S_{\mathrm{even}}/S_{\mathrm{odd}} = 2/3\) получаем \(\delta N \approx 0.1333\). Следовательно, эффективное число калибровочных степеней свободы на лабораторных масштабах равно
	\begin{equation}
		N_{\mathrm{eff}} = 12 + \delta N = 12 + \beta\gamma\,\frac{S_{\mathrm{even}}}{S_{\mathrm{odd}}}
		\approx 12.1333 .
	\end{equation}
	Предсказываемая обратная постоянная тонкой структуры есть
	\begin{equation}
		\boxed{\alpha^{-1}_{\mathrm{YUCT}} =
			\left(\frac{S_{\mathrm{odd}}}{S_{\mathrm{even}}}\right)^{\!N_{\mathrm{eff}}}
			= \left(\frac{3}{2}\right)^{\!12 + \beta\gamma\,S_{\mathrm{even}}/S_{\mathrm{odd}}} } .
		\label{eq:alphaYUCT2}
	\end{equation}
	Численно,
	\[
	\alpha^{-1}_{\mathrm{YUCT}} = \left(\frac{3}{2}\right)^{\!12.1333} \approx 137.0 ,
	\]
	что отличается от значения CODATA 2022 \(\alpha^{-1}_{\mathrm{exp}} = 137.035999084\) менее чем на \(0.03\%\). Уравнение (\ref{eq:alphaYUCT2}) не содержит свободных параметров: \(12\) — топологический инвариант 19-мерного координационного пространства, \(S_{\mathrm{odd}}/S_{\mathrm{even}}\) и \(\beta\gamma\) — фундаментальные константы YUCT, уже зафиксированные независимыми наблюдениями (Приложения Ferra1.2 и P).
	
	\subsubsection{Слабая связь}
	Слабое взаимодействие не требует отдельного эффективного числа мод. Угол Вайнберга \(\theta_W\) задаётся отношением масс \(m_W/m_Z = \cos\theta_W\), и обе массы \(m_W\) и \(m_Z\) выводятся из глубины \(L_W = 96\) и соответствующих резонансных факторов (Приложение \(\Lambda\)). Таким образом, слабая связь \(\alpha_W = \alpha/\sin^2\theta_W\) определяется без дополнительных параметров.
	
	\subsubsection{Сильная связь}
	Квантовая хромодинамика имеет 8 генераторов глюонов, поэтому обратная связь на планковском масштабе равна
	\begin{equation}
		\alpha_s^{-1}(M_{\mathrm{Pl}}) = \left(\frac{3}{2}\right)^{\!8 + \beta\gamma\,S_{\mathrm{even}}/S_{\mathrm{odd}}}
		\approx 25.5 \qquad (\alpha_s \approx 0.039) .
	\end{equation}
	Понижение энергетического масштаба в YUCT описывается как последовательная активация дополнительных координационных оболочек, каждая из которых даёт дискретный вклад \(\Delta L\) в обратную связь. В континуальном пределе это воспроизводит логарифмический бег QCD, но с естественной дискретизацией на глубинах, кратных \(S_{\mathrm{odd}}\) и \(S_{\mathrm{even}}\).
	
	\subsubsection{Юкавские связи}
	Юкавская связь заряженного фермиона \(f\) равна
	\begin{equation}
		y_f = \frac{\sqrt{2}\,m_f}{v},
		\qquad
		m_f = M_{\mathrm{Pl}}\left(\frac{2}{3}\right)^{\!96+k_f}\xi_f,
		\qquad
		v   = M_{\mathrm{Pl}}\left(\frac{2}{3}\right)^{\!96}\xi_W .
	\end{equation}
	Следовательно,
	\begin{equation}
		y_f = \sqrt{2}\,\frac{\xi_f}{\xi_W}\,\left(\frac{2}{3}\right)^{\!k_f} .
	\end{equation}
	Резонансный фактор \(\xi_f\) выражается через функцию совместимости \(C_{fH}(L_f,L_H)\), введённую в «Координационной систематике YUCT», с универсальной шириной \(\sigma \approx 0.20\). Таким образом, все юкавские связи фиксируются известными координационными глубинами фермионов и бозона Хиггса.
	
	\subsubsection{Дискретная лестница эффективных мод}
	Рисунок~\ref{fig:ladder} иллюстрирует, как эффективное число калибровочных мод \(N_{\mathrm{eff}}\) меняется с глубиной \(L\) (уменьшением энергии). Каждая ступень соответствует порогу массы, где частица (или семейство частиц) перестаёт вносить вклад в вакуумную поляризацию. Гладкая огибающая — это низкоэнергетический предел, определяющий наблюдаемую константу связи.
	
	\begin{figure}[htbp]
		\centering
		\begin{tikzpicture}[scale=0.9]
			% оси
			\draw[->] (0,0) -- (16,0) node[right] {\(L\) (глубина)};
			\draw[->] (0,0) -- (0,8) node[above] {\(N_{\mathrm{eff}}\)};
			% ступени
			\draw[thick,blue] (0, 5.0) -- (1.5, 5.0);
			\draw[thick,blue] (1.5, 5.5) -- (3.2, 5.5);
			\draw[thick,blue] (3.2, 6.0) -- (5.0, 6.0);
			\draw[thick,blue] (5.0, 6.5) -- (7.0, 6.5);
			\draw[thick,blue] (7.0, 7.0) -- (9.0, 7.0);
			\draw[thick,blue] (9.0, 7.5) -- (11.0, 7.5);
			\draw[thick,blue] (11.0, 7.8) -- (13.0, 7.8);
			\draw[thick,blue] (13.0, 7.9) -- (15.0, 7.9) node[right,blue] {12.1333};
			% подписи
			\node[below] at (1.5,0) {$L_W$ (96)};
			\node[below] at (3.2,0) {$L_Z$};
			\node[below] at (5.0,0) {$L_{\mathrm{top}}$};
			\node[below] at (7.0,0) {$L_{\mathrm{Higgs}}$};
			\node[below] at (9.0,0) {$L_{\tau}$};
			\node[below] at (11.0,0) {$L_{\mu}$};
			\node[below] at (13.0,0) {$L_e$ (107)};
			% гладкая огибающая
			\draw[red, dashed, thick] plot[smooth] coordinates {(0,5.0) (1.5,5.5) (3.2,6.0) (5.0,6.5) (7.0,7.0) (9.0,7.5) (11.,7.8) (13.,7.9)};
		\end{tikzpicture}
		\caption{Эффективное число калибровочных мод \(N_{\mathrm{eff}}\) как функция координационной глубины \(L\). Ступени соответствуют массовым порогам, пунктирная кривая — низкоэнергетический предел, определяющий постоянную тонкой структуры.}
		\label{fig:ladder}
	\end{figure}
	
	С этими результатами все безразмерные параметры Стандартной модели сводятся к топологическому инварианту 12, фундаментальному отношению \(S_{\mathrm{odd}}/S_{\mathrm{even}}\), универсальной поправке \(\beta\gamma\) и координационным глубинам частиц. Свободных непрерывных параметров не остаётся.
	
	\subsubsection{Полное устранение свободных непрерывных параметров в Стандартной модели}
	
	Полученные выше результаты показывают, что при встраивании в рамки YUCT в стандартной квантовой теории поля \emph{не остаётся свободных непрерывных параметров}. Таблица~\ref{tab:SMparams} суммирует, как каждая независимая константа Стандартной модели фиксируется фундаментальными числами координационной теории вместе с координационными глубинами соответствующих частиц.
	
	\begin{table}[htbp]
		\centering
		\caption{Фиксация безразмерных параметров Стандартной модели в YUCT}
		\label{tab:SMparams}
		\small
		\begin{tabular}{@{}p{3.0cm} p{5.2cm} p{5.2cm}@{}}
			\toprule
			\textbf{Параметр} & \textbf{Выражение в YUCT} & \textbf{Ключевые ссылки} \\
			\midrule
			Постоянная тонкой структуры \(\alpha\) &
			\(\alpha^{-1} = (S_{\mathrm{odd}}/S_{\mathrm{even}})^{12+\beta\gamma\,S_{\mathrm{even}}/S_{\mathrm{odd}}}\) &
			Данная работа, Приложение P, Ferra1.2 \\
			Угол Вайнберга \(\theta_W\) &
			Выводится из \(m_W/m_Z\), обе массы следуют из глубины \(L_W=96\) &
			Приложение \(\Lambda\) \\
			Сильная связь \(\alpha_s\) &
			\(\alpha_s^{-1} = (S_{\mathrm{odd}}/S_{\mathrm{even}})^{8+\beta\gamma\,S_{\mathrm{even}}/S_{\mathrm{odd}}}\) на планковском масштабе; бег получается из дискретных оболочек глубины &
			Данная работа, Приложение P \\
			Юкавские связи \(y_f\) &
			\(y_f = \sqrt{2}\,(\xi_f/\xi_W)\,(S_{\mathrm{even}}/S_{\mathrm{odd}})^{k_f}\) с резонансными факторами \(\xi_f,\xi_W\) из матрицы совместимости &
			Приложение J, Координационная систематика \\
			Массы фермионов & Полуцелое квантование \(m_f = m_e\,q^{N_f}\), \(q=(S_{\mathrm{odd}}/S_{\mathrm{even}})^{1/3}\) &
			Приложение J, Координационная систематика \\
			CP-нарушающий \(\theta\)-параметр &
			\(\theta \sim (S_{\mathrm{even}}/S_{\mathrm{odd}})^{L_W/2 + S_{\mathrm{odd}}/S_{\mathrm{even}}}\) &
			«Решение сильной CP-проблемы» \\
			Космологическая постоянная \(\Omega_\Lambda\) &
			\(\Omega_\Lambda = 12/18 = 2/3\) (топологический инвариант) &
			Данная работа, Приложение A \\
			\bottomrule
		\end{tabular}
	\end{table}
	
	Все записи в таблице опираются лишь на три целых числа, характеризующих координационную сеть — базовую глубину \(L_W=96\), число калибровочных мод \(N_{\mathrm{gauge}}=12\) и число глюонных мод \(N_{\mathrm{gluon}}=8\) — вместе с универсальным отношением \(S_{\mathrm{odd}}/S_{\mathrm{even}}\), произведением \(\beta\gamma\) и дискретными координационными глубинами \(L_f\) фермионов и бозона Хиггса. Резонансный фактор \(\xi_f/\xi_W\) не является подгоночным параметром; он вычисляется из матрицы совместимости \(C_{fH}\) с эмпирически устойчивой шириной \(\sigma \approx 0.20\) (см. Координационную систематику YUCT). CP-нечётный \(\theta\)-параметр аналогичным образом фиксируется глубиной электрослабого масштаба и фундаментальным смещением порядок–хаос. Следовательно, весь набор безразмерных констант, появляющихся в Стандартной модели, определяется без каких-либо произвольных непрерывных входных данных.
	
	Это структурное сокращение завершает программу превращения квантовой теории поля из теории с 19 свободными параметрами в эффективное описание, числовое содержание которого полностью диктуется координационной архитектурой YUCT.
	
	\section{Единое координационное поле и два протокола}
	\label{sec:unified-protocols}
	
	Изложенная выше модель геометрического натяжения может быть сделана более точной путём вложения в полную 19-мерную геометрию YUCT. Как показано ранее, космологическая постоянная \(\Omega_\Lambda=2/3\) возникает из топологии компактного слоя \(F_{18}\). Более глубокий анализ показывает, что то же координационное поле \(\Psi_{MN}\), порождающее \(\Theta_{\mu\nu}\), даёт также два фундаментальных координационных протокола: централизованный YPSDC и децентрализованный d‑YPSDC (полный вывод см. в статье о первых принципах). Единое действие имеет вид
	\begin{equation}
		S_{19} = \frac{1}{2\kappa_{19}} \int d^{19}X \sqrt{-G}\, R(G)
		+ \int d^{19}X \sqrt{-G}\,
		\Bigl[ -\frac14 \Tr(\Psi_{MN}\Psi^{MN})
		+ \frac12 G^{MN}\partial_M\phi \partial_N\phi
		- V(\phi) \Bigr],
		\label{eq:unified-gravity}
	\end{equation}
	где \(\phi = \ln(K_{\mathrm{eff}}/K_0)\) — эффективное 4-мерное скалярное поле, использованное в этой работе. После компактификации на \(F_{18}\) нулевая мода \(\Psi_{MN}\) даёт в точности скаляр \(\phi\) и тензор геометрического натяжения \(\Theta_{\mu\nu}\) из уравнения (\ref{eq:Theta}).
	
	Два протокола соответствуют различным динамическим режимам одного и того же поля, различающимся источником индекса \(\kappa\):
	\begin{enumerate}[leftmargin=*]
		\item \textbf{Централизованный YPSDC.} Один агент активно возбуждает \(\Psi_{MN}\), создавая локализованное возмущение, распространяющееся на другого агента. Отправитель действует как точечный источник \(D_M\Psi^{MN} = J^N_{\rm sender}\). Это знакомая картина коммуникации, которая в гравитационном контексте соответствует локальному переуплотнению, порождающему ньютоновский потенциал.
		\item \textbf{Децентрализованный d‑YPSDC.} Все агенты пассивно считывают один и тот же индекс среды из длинноволновой фоновой моды \(\Psi_{MN}\). Отправитель не нужен; корреляции обеспечиваются глобальной топологией \(F_{18}\) и общим словарём. В гравитации этот режим даёт космологическую постоянную, тёмную энергию и нелокальные корреляции, наблюдаемые в квантовой запутанности.
	\end{enumerate}
	
	Переход между режимами контролируется безразмерным параметром \(\Theta = |J_{\rm agent}| / |J_{\rm env}|\), где \(J_{\rm agent}\) — ток активного отправителя, а \(J_{\rm env}\) — эффективный ток космологического фона. При \(\Theta\ll1\) (децентрализованный режим) поле \(\phi\) почти однородно, \(\partial_M\phi\approx0\), и тензор геометрического натяжения сводится к эффективной космологической постоянной плюс малая компонента тёмной материи. При \(\Theta\gg1\) (централизованный режим) доминируют локальные градиенты \(\phi\), давая ньютоновский потенциал и поправку \(1/r\), объясняющую плоские галактические кривые вращения.
	
	Это объединение имеет глубокие следствия для квантовой физики. Квантовая запутанность больше не является «жутким дальнодействием», а просто предельным случаем децентрализованного d‑YPSDC: запутанная пара разделяет общий словарь (волновую функцию), и измерение на одной частице даёт индекс, который поле \(\Psi_{MN}\) одновременно доставляет другой частице, без какого-либо сверхсветового сигнала. Степень нарушения неравенств Белла связана с эффективностью координации:
	\begin{equation}
		S = 2\sqrt{2}\,\frac{K_{\mathrm{eff}}}{K_{\mathrm{eff}}+1},
		\label{eq:bell-eff}
	\end{equation}
	так что стандартный квантовый максимум \(2\sqrt{2}\) достигается только при \(K_{\mathrm{eff}}\to\infty\) (идеальный словарь), тогда как при конечной эффективности корреляции слабее.
	
	Таким образом, модель гравитации как геометрического натяжения, космологическая постоянная, иерархия масс фермионов и квантовая нелокальность следуют из одной сущности — координационного поля \(\Psi_{MN}\) на \(\mathcal{M}_{19}=M_4\times F_{18}\). Два протокола не являются отдельными аксиомами, а представляют собой выведенные динамические фазы одной и той же фундаментальной физики. Количественный анализ фазового перехода между двумя режимами и его наблюдательных проявлений будет представлен в другом месте.
	
	\section{Циклическая космология}
	
	Потенциал \(V(\phi)\) имеет единственный глобальный минимум при \(\phi=0\). В процессе гравитационного коллапса скалярное поле смещается в сторону больших положительных значений. Когда \(\phi\) превышает критический порог \(\phi_{\mathrm{crit}}\), эффективное давление становится достаточно отрицательным, чтобы вызвать тахионную неустойчивость. Условие имеет вид
	\[
	V''(\phi) > \frac{9}{4} H^2.
	\]
	Используя уравнение Фридмана–Якушева
	\[
	H^2 = \frac{8\pi G_0}{3} \left[ \rho_b + \frac{\kappa}{2}\dot\phi^2 + V(\phi) \right],
	\]
	и замечая, что вблизи отскока \(\dot\phi^2\) велико, находим, что неустойчивость возникает, когда
	\[
	e^{\lambda\phi_{\mathrm{crit}}} \approx \frac{9}{4} \frac{\kappa\dot\phi^2}{\lambda^2 V_0} + \cdots
	\]
	Детальное численное интегрирование показывает \(\phi_{\mathrm{crit}} \approx \ln\Phi\), где \(\Phi \approx 1.618\) — золотое сечение. Затем поле быстро скатывается к \(\phi=0\), высвобождая накопленную энтропию и вызывая короткий период экспоненциального расширения (инфляцию). Цикл повторяется бесконечно, избегая начальной сингулярности.
	
	\section{Крупномасштабная структура космоса и природа фундаментальных масштабов}
	\label{sec:sponge}
	
	Уравнения координационного поля \(\phi\) не только определяют локальную гравитационную силу, но и формируют глобальную архитектуру Вселенной. В этом разделе мы показываем, как «губкоподобная» космическая сеть, возможность причинно несвязанных дочерних вселенных и не универсальность планковской длины возникают как естественные следствия модели геометрического натяжения.
	
	\subsection{Губкоподобная морфология космической сети}
	В эпоху доминирования материи уравнение поля~\eqref{eq:phieom} допускает стационарные нелинейные решения, разделяющие области с разной координационной глубиной \(L\). Профиль одномерной плоской стенки между двумя областями \(\phi_1\) и \(\phi_2\) имеет вид
	\begin{equation}
		\phi(x) = \frac{2}{\lambda}\,\ln\!\left[
		\frac{1 + \tanh(\lambda\sqrt{V_0/\kappa}\, x/2)}
		{1 - \tanh(\lambda\sqrt{V_0/\kappa}\, x/2)}
		\right],
		\label{eq:domain-wall}
	\end{equation}
	гладко интерполируя между двумя вакуумами. Поверхностное натяжение такой стенки равно \(\sigma \simeq \sqrt{\kappa V_0}/\lambda^2\).
	
	В трёх измерениях стенки образуют связную ячеистую сеть: узлы — глубокие кластеры (\(\phi\) велико), нити — пересечения стенок, а пустоты — области, где \(\phi \approx 0\). Это в точности «космическая губка», наблюдаемая в обзорах галактик. Гравитационно-линзовый сигнал стенки усиливается градиентом \(\phi\), давая эффективную поверхностную плотность тёмной материи:
	\begin{equation}
		\Sigma_{\rm DM}^{\rm eff} = \frac{\kappa}{4\pi G_0} \int_{-\infty}^{+\infty} \!dz\; \nabla^2\delta\phi(z)
		\simeq \frac{\sqrt{\kappa V_0}}{2\pi G_0\,\lambda},
		\label{eq:wall-lensing}
	\end{equation}
	что может быть проверено с помощью суммированных данных слабого линзирования.
	
	\subsection{Отключённые дочерние вселенные и «неопределённое расстояние»}
	Если локальная флуктуация \(\phi\) превышает критическое значение \(\phi_{\rm crit}=\ln\Phi\), происходит \textbf{локальный} фазовый переход: маленький пузырь нового вакуума нуклеируется и расширяется, создавая причинно несвязную область — дочернюю вселенную. Две вселенные разделены доменной стенкой с поверхностным слоем очень крутого \(\nabla\phi\). Расстояние между ними не может быть измерено никакой классической геодезической, потому что аффинный параметр вдоль кривой, пересекающей стенку, расходится. Единственное осмысленное «расстояние» — это \textbf{разность координационных глубин}
	\begin{equation}
		\Delta L = L_{\rm parent} - L_{\rm child},
		\label{eq:deltaL}
	\end{equation}
	которая является топологическим квантовым числом. Следовательно, «неопределённое расстояние» между вселенными есть прямое следствие дискретной структуры координационных слоёв.
	
	\subsection{Планковская длина как локальный, зависящий от окружения масштаб}
	В YUCT фундаментальной длиной является комптоновская длина волны самого лёгкого координационного кластера с глубиной \(L\). Используя формулу массы \(m \simeq M_{\rm Pl}\,(2/3)^L\), эта длина равна
	\begin{equation}
		\ell_{\rm coord}(L) = \frac{\hbar}{m c}
		= \frac{\hbar}{M_{\rm Pl}c} \left(\frac{3}{2}\right)^{\!L}
		= \ell_{\rm Pl} \left(\frac{3}{2}\right)^{\!L}.
		\label{eq:length-from-L}
	\end{equation}
	Для вакуума текущей эпохи \(L_{\rm vacuum}\to\infty\), но любой конечный кластер зондирует масштаб \(\ell_{\rm coord}(L) \gg \ell_{\rm Pl}\). Обычная планковская длина — это просто частный случай этой общей формулы при \(L=1\), т.е. она характеризует координационный слой, доминирующий в сегодняшних низкоэнергетических явлениях.
	
	Следовательно, гравитационная постоянная также зависит от слоя:
	\begin{equation}
		G_{\rm eff}(\phi) = \frac{G_0}{1-4\pi G_0\alpha_2\phi^2},
		\label{eq:Geff-layer}
	\end{equation}
	так что в областях с большим \(\phi\) гравитация слабее. Это даёт естественное разрешение напряжения Хаббла: локальные измерения скорости расширения зондируют немного другую среду \(\phi\), чем реликтовое излучение, что приводит к систематически отличающемуся значению \(H_0\). Количественное исследование этого эффекта будет представлено в другом месте.
	
	\section{Проверяемые предсказания}
	
	\subsection{Галактические кривые вращения}
	Член модифицированной гравитации предсказывает плоские кривые вращения с нормировкой, определяемой космической долей тёмной материи. Детальное согласование с индивидуальными галактиками будет представлено в отдельной работе.
	
	\subsection{Рост космических структур}
	Координационное поле модифицирует линейный фактор роста. Используя стандартную теорию возмущений, дробное отклонение от \(\Lambda\)CDM на красном смещении \(z=0.5\) составляет примерно \(2\)–\(3\%\), что измеримо обзорами Euclid и DESI.
	
	\subsection{Ограничения на космологические параметры}
	Модель предсказывает \(\Omega_\Lambda = 2/3\) точно и \(\Omega_{\mathrm{DM}}^{\mathrm{eff}} = 1 - 0.05 - 2/3 \approx 0.2833\). Эти значения могут быть сопоставлены с данными Planck, DES и LSST.
	
	\section{Заключение}
	Мы вывели гравитацию из информационно-теоретических аксиом YUCT. Тензор геометрического натяжения \(\Theta_{\mu\nu}\) объединяет тёмную материю и тёмную энергию, космологическая постоянная фиксируется топологией, а циклическая космическая история возникает естественным образом. Все параметры выражаются через планковский масштаб и константы словаря. Теория проверяема с помощью современных космологических обзоров.
	
	\begin{thebibliography}{99}
		\bibitem{Y-appA} A. V. Yakushev, YUCT Lagrangian V35.0 (Приложение A).
		\bibitem{Y-appL} A. V. Yakushev, Fractal Coordination Error Scaling (Приложение L).
		\bibitem{Y-appM} A. V. Yakushev, Cosmology -- Inflation as a Coordination Phase Transition (Приложение M).
		\bibitem{Y-appG} A. V. Yakushev, Resolution of the Quantum Gravity Problem (Приложение G).
	\end{thebibliography}
	
\end{document}